\chapter{Сходимост и~закони за големите числа}

\begin{rembox}[Забележка]
Това е сред най-важните теми в~курса - сходимостта на случайни
величини и~законите за големите числа стоят в~основата на цялата
математическа статистика (оценки, доверителни интервали, гранични
теореми).
\end{rembox}

\section{Видове сходимост}

\begin{rembox}[Постановка]
Нека $(X_n)_{n=1}^{\infty}$ е редица от случайни величини във
вероятностно пространство $V=(\Omega,\mathcal A,\mathbb P)$.
\end{rembox}

\begin{defbox}[~(Почти сигурна сходимост, п.с.)]
Нека $L_X = \{\omega\in\Omega : \lim_{n\to\infty}X_n(\omega)=X(\omega)\}$.
Казваме, че $X_n\xrightarrow[n\to\infty]{\text{п.с.}} X$, ако
$\mathbb P(L_X)=1$ (,,$X_n$ клони почти сигурно към $X$``). Възможно
е да съществуват изходи $\omega$, за които границата не е изпълнена,
но множеството от такива изходи да има вероятност нула.
\end{defbox}

\begin{defbox}[~(Сходимост по вероятност)]
Нека $A_{n,\varepsilon} = \{\omega\in\Omega : |X_n(\omega)-X(\omega)|
>\varepsilon\}$. Казваме, че $X_n\xrightarrow[n\to\infty]{\mathbb P}
X$, ако $\forall\varepsilon>0$: $\lim_{n\to\infty}\mathbb
P(A_{n,\varepsilon})=0$.
\end{defbox}

\begin{defbox}[~(Сходимост по разпределение)]
Нека $C_{F_X}=\{x\in\mathbb R : F_X\text{ е непрекъсната в~} x\}$ -
множеството от точки на непрекъснатост на функцията на разпределение
$F_X$. Казваме, че $X_n\xrightarrow[n\to\infty]{d} X$, ако за всяко
$x\in C_{F_X}$: $\lim_{n\to\infty}F_{X_n}(x)=F_X(x)$.
\end{defbox}

\begin{thmbox}[]{Теорема}
Нека $(X_n)_{n=1}^{\infty}$ е редица от случайни величини.
\begin{enumerate}[label=\alph*)]
  \item Ако $X_n\xrightarrow{\text{п.с.}}X$, то $X_n\xrightarrow{\mathbb P}X$.
  \item Ако $X_n\xrightarrow{\mathbb P}X$, то $X_n\xrightarrow{d}X$.
\end{enumerate}
(Т.е.\ п.с.\ сходимостта е най-силна, следвана от сходимостта по
вероятност, а~сходимостта по разпределение е най-слабата.)
\end{thmbox}

\begin{proofbox}
\textbf{а)} Нека $X_n\xrightarrow{\text{п.с.}}X$, т.е.\ $\mathbb
P(L)=1$ за $L=\{\omega:\lim_n X_n(\omega)=X(\omega)\}$. Твърдим, че
\[
  L = \bigcap_{r=1}^{\infty}\bigcup_{n=1}^{\infty}\bigcap_{k=n}^{\infty} \Big\{\omega\in\Omega : |X_k(\omega)-X(\omega)|\le\frac1r\Big\}.
\]
\emph{($\subseteq$)} Нека $\tilde\omega\in L$. За всяко $r\ge1$
(с~$\varepsilon=1/r$), понеже $X_n(\tilde\omega)\to X(\tilde\omega)$,
съществува $n_r$ такова, че $|X_n(\tilde\omega)-X(\tilde\omega)|<
\tfrac1r$ за $n\ge n_r$, т.е.\ $\tilde\omega\in\bigcap_{k=n_r}^{\infty}
\{|X_k-X|\le\tfrac1r\}\subseteq\bigcup_n\bigcap_{k=n}^{\infty}\{\dots\}$,
и~това важи за всяко $r$, следователно $\tilde\omega$ принадлежи на
дясната страна.

\emph{($\supseteq$)} Нека $\tilde\omega$ принадлежи на дясната
страна. За всяко $r\ge1$ съществува $n_r$, такова че $\tilde\omega\in
\bigcap_{k=n_r}^{\infty}\{|X_k-X|\le\tfrac1r\}$, т.е.\
$|X_k(\tilde\omega)-X(\tilde\omega)|\le\tfrac1r$ за всяко $k\ge n_r$.
Това означава $\lim_{k\to\infty}X_k(\tilde\omega)=X(\tilde\omega)$
(по дефиниция на граница), т.е.\ $\tilde\omega\in L$.

От равенството и~закона на Де Морган:
\[
  L^c = \bigcup_{r=1}^{\infty}\bigcap_{n=1}^{\infty}\underbrace{\bigcup_{k=n}^{\infty}A_{k,\frac1r}}_{=:B_{n,r}}.
\]
Понеже $\mathbb P(L)=1$, имаме $\mathbb P(L^c)=0$, следователно
$\mathbb P\big(\bigcap_n B_{n,r}\big)=0$ за всяко $r\ge1$ (обединение
на неотрицателни членове е нула точно когато всеки е нула). Нека
$r\ge1$ е фиксирано. Редицата $(B_{n,r})_{n\ge1}$ намалява по $n$
(защото обединяваме над все по-малко индекси), следователно по
непрекъснатостта на вероятността отгоре (Тема~2):
\[
  0 = \mathbb P\Big(\bigcap_{n=1}^{\infty}B_{n,r}\Big) = \lim_{n\to\infty}\mathbb P(B_{n,r}).
\]
Понеже $B_{n,r}=\bigcup_{k=n}^{\infty}A_{k,1/r}\supseteq A_{n,1/r}$,
имаме $\mathbb P(A_{n,1/r})\le\mathbb P(B_{n,r})\to0$, т.е.
$\lim_{n\to\infty}\mathbb P(A_{n,1/r})=0$ за всяко $r\ge1$ - а~това
означава точно $X_n\xrightarrow{\mathbb P}X$ (произволно $\varepsilon>0$
може да бъде „уловено“ от достатъчно голямо $r$ с~$1/r<\varepsilon$).

\textbf{б)} Нека $X_n\xrightarrow{\mathbb P}X$, т.е.\
$\lim_n\mathbb P(A_{n,\varepsilon})=0$ за всяко $\varepsilon>0$. Нека
$x\in C_{F_X}$ е фиксирано и~нека $\varepsilon>0$. Забелязваме
двете включвания:
\[
  \{X_n<x\} \subseteq A_{n,\varepsilon}\cup\{X<x+\varepsilon\},
  \qquad
  \{X<x-\varepsilon\}\cap A_{n,\varepsilon}^c \subseteq \{X_n<x\}
\]
(геометрично: ако $X_n<x$ и~$|X_n-X|\le\varepsilon$, то $X<x+
\varepsilon$; а~ако $X<x-\varepsilon$ и~$|X_n-X|\le\varepsilon$, то
$X_n<x$). От тези включвания следва:
\[
  \mathbb P(X<x-\varepsilon) - \mathbb P(A_{n,\varepsilon}) \le \mathbb P(X_n<x) \le \mathbb P(X<x+\varepsilon)+\mathbb P(A_{n,\varepsilon}),
\]
т.е.\ $F_X(x-\varepsilon)-\mathbb P(A_{n,\varepsilon})\le
\mathbb P(X_n<x)\le F_X(x+\varepsilon)+\mathbb P(A_{n,\varepsilon})$.
Взимаме $\liminf$/$\limsup$ по $n\to\infty$ (при което $\mathbb
P(A_{n,\varepsilon})\to0$):
\[
  F_X(x-\varepsilon) \le \liminf_{n\to\infty}\mathbb P(X_n<x) \le \limsup_{n\to\infty}\mathbb P(X_n<x) \le F_X(x+\varepsilon).
\]
Понеже $x\in C_{F_X}$ (точка на непрекъснатост), $\lim_{\varepsilon
\to0}F_X(x-\varepsilon)=\lim_{\varepsilon\to0}F_X(x+\varepsilon)=
F_X(x)$, следователно при $\varepsilon\to0$ границите отляво и~отдясно се сближават до $F_X(x)$, откъдето
$\lim_{n\to\infty}\mathbb P(X_n<x)=F_X(x)$, т.е.\ $X_n
\xrightarrow{d}X$. \qedbg
\end{proofbox}

\begin{thmbox}[]{Теорема}
Нека $(X_n)_{n=1}^{\infty}$ е редица от случайни величини в~едно
вероятностно пространство. Ако $X_n\xrightarrow{d}c$, то
$X_n\xrightarrow{\mathbb P}c$, където $c$ е константа.
\end{thmbox}

\begin{proofbox}
$F_c(x) = I_{\{x>c\}}$, т.е.\ $C_{F_c}=\mathbb R\setminus\{c\}$.
Фиксираме $\varepsilon>0$ и~разписваме:
\begin{align*}
  \mathbb P(|X_n-c|>\varepsilon) &= \mathbb P(X_n>c+\varepsilon)+\mathbb P(X_n<c-\varepsilon)\\
  &= 1-\mathbb P(X_n\le c+\varepsilon) + F_{X_n}(c-\varepsilon)\\
  &= 1 - F_{X_n}(c+\varepsilon) - \mathbb P(X_n=c+\varepsilon) + F_{X_n}(c-\varepsilon).
\end{align*}
Понеже $c+\varepsilon,c-\varepsilon\in C_{F_c}$, от $X_n
\xrightarrow{d}c$ следва $F_{X_n}(c+\varepsilon)\to F_c(c+\varepsilon)
=1$ и~$F_{X_n}(c-\varepsilon)\to F_c(c-\varepsilon)=0$ при
$n\to\infty$. Следователно
\[
  0 \le \lim_{n\to\infty}\mathbb P(|X_n-c|>\varepsilon) = \underbrace{1-1-0}_{=0} - \lim_{n\to\infty}\mathbb P(X_n=c+\varepsilon) = -\lim_{n\to\infty}\mathbb P(X_n=c+\varepsilon) \le 0
\]
(последното неравенство, защото вероятности са неотрицателни), т.е.\
границата е едновременно $\ge0$ и~$\le0$, откъдето $\lim_{n\to
\infty}\mathbb P(|X_n-c|>\varepsilon)=0$. \qedbg
\end{proofbox}

\section{Неравенство на Чебишов (обобщение)}

\begin{thmbox}[~(на Чебишов, чрез индикатор)]{Теорема}
Нека $X$ е случайна величина с~$DX<\infty$. Тогава за всяко $a>0$:
\[
  \mathbb P(|X-\mathbb EX|>a) \le \frac{DX}{a^2}.
\]
\end{thmbox}

\begin{proofbox}
Нека $Y=X-\mathbb EX$. Тогава $\mathbb P(|Y|>a)=\mathbb E\mathbb
1_{\{|Y|>a\}}$, и~разбивайки по допълващи се събития:
\begin{align*}
  DX = \mathbb EY^2 &= \mathbb E\big[Y^2\cdot1\big] = \mathbb E\big[Y^2\big(I_{\{|Y|\le a\}}+I_{\{|Y|>a\}}\big)\big]\\
  &\ge \mathbb E\big[Y^2I_{\{|Y|>a\}}\big] \ge a^2\,\mathbb EI_{\{|Y|>a\}} = a^2\,\mathbb P(|Y|>a)
\end{align*}
(изпуснахме неотрицателния член с~$I_{\{|Y|\le a\}}$, а~на
събитието $\{|Y|>a\}$ имаме $Y^2>a^2$). Следователно $\mathbb
P(|X-\mathbb EX|>a)\le DX/a^2$. \qedbg
\end{proofbox}

\begin{corbox}[~(обобщение за по-висок момент)]
Нека $X$ е случайна величина и~$n\ge1$. Тогава за всяко $a>0$:
\[
  \mathbb P(|X-\mathbb EX|>a) \le \frac{\mathbb E|X-\mathbb EX|^n}{a^n}
\]
(доказва се по абсолютно същата схема, замествайки квадрат с~$n$-та степен: на $\{|Y|>a\}$ имаме $|Y|^n>a^n$).
\end{corbox}

\section{Закони за големите числа}

\begin{defbox}[~(Независими еднакво разпределени - i.i.d.)]
Редицата $(X_i)_{i=1}^{n}$ (или $(X_i)_{i=1}^{\infty}$) е от
\textbf{независими еднакво разпределени} случайни величини
(съкратено \textbf{i.i.d.}), ако $X_1\stackrel{d}{=}X_i$ (т.е.\
$F_{X_1}=F_{X_i}$) за всяко $i\ge2$, и~$(X_i)$ са независими в~съвкупност (т.е.\ за всяко $n\ge1$ и~всеки избор на индекси
$(i_1,\dots,i_n)$, величините $(X_{i_1},\dots,X_{i_n})$ са
независими в~съвкупност).
\end{defbox}

\begin{defbox}[~(Слаб закон за големите числа, СЗГЧ)]
Нека $(X_i)_{i=1}^{\infty}$ е редица от случайни величини с~очаквания $(\mathbb EX_i)_{i=1}^{\infty}$. Казваме, че за редицата е
в~сила \textbf{слабият закон за големите числа}, ако
\[
  \frac{\sum_{i=1}^n (X_i-\mathbb EX_i)}{n} \xrightarrow[n\to\infty]{\mathbb P} 0.
\]
\end{defbox}

\begin{thmbox}[~(СЗГЧ)]{Теорема}
Нека $(X_i)_{i=1}^{\infty}$ е i.i.d.\ редица с~$\mathbb EX_1$ и~$DX_1<\infty$. Тогава
\[
  \frac{\sum_{i=1}^n X_i}{n} \xrightarrow[n\to\infty]{\mathbb P} \mathbb EX_1.
\]
\end{thmbox}

\begin{proofbox}
Фиксираме $\varepsilon>0$. Понеже $(X_i)$ са i.i.d., $D\big(
\sum_{i=1}^n(X_i-\mathbb EX_i)\big) = \sum_{i=1}^n D(X_i-\mathbb
EX_i) = \sum_{i=1}^n DX_i = n\,DX_1$ (адитивност на дисперсията за
независими събираеми, Тема~5/9; изваждането на константа $\mathbb
EX_i$ не променя дисперсията). По неравенството на Чебишов,
приложено към $\sum_{i=1}^n(X_i-\mathbb EX_i)$:
\begin{align*}
  \mathbb P\Big(\Big|\frac{\sum_{i=1}^n X_i}{n}-\mathbb EX_1\Big|>\varepsilon\Big) &= \mathbb P\Big(\Big|\sum_{i=1}^n(X_i-\mathbb EX_i)\Big|>n\varepsilon\Big) \le \frac{D\big(\sum_{i=1}^n(X_i-\mathbb EX_i)\big)}{\varepsilon^2n^2}\\
  &= \frac{n\,DX_1}{\varepsilon^2n^2} = \frac{DX_1}{n\varepsilon^2} \xrightarrow[n\to\infty]{} 0. \qedbg
\end{align*}
\end{proofbox}

\begin{defbox}[~(Усилен закон за големите числа, УЗГЧ)]
Нека $(X_i)_{i=1}^{\infty}$ е редица от случайни величини с~очаквания $(\mathbb EX_i)_{i=1}^{\infty}$. Казваме, че за редицата е
в~сила \textbf{усиленият закон за големите числа}, ако
\[
  \frac{\sum_{i=1}^n (X_i-\mathbb EX_i)}{n} \xrightarrow[n\to\infty]{\text{п.с.}} 0.
\]
\end{defbox}

\begin{thmbox}[~(УЗГЧ)]{Теорема}
Нека $(X_i)_{i=1}^{\infty}$ е i.i.d.\ редица с~$\mathbb E|X_1|<\infty$.
Тогава
\[
  \frac{\sum_{i=1}^n X_i}{n} \xrightarrow[n\to\infty]{\text{п.с.}} \mathbb EX_1.
\]
\end{thmbox}

\begin{proofbox}
\emph{(Доказваме при допълнителното техническо предположение
$\mathbb EX_1^4<\infty$, което опростява сметките.)} Без ограничение
на общността приемаме $\mathbb EX_1=0$ (иначе разглеждаме
$X_i-\mathbb EX_i$). Полагаме $S_k=\sum_{i=1}^k X_i$. Целта е да
докажем $\mathbb P(L)=1$ за $L=\{\omega: S_n(\omega)/n\to0\}$, т.е.\
$\mathbb P(L^c)=0$. Аналогично на предната теорема,
\[
  L^c = \bigcup_{r=1}^{\infty} C_r, \qquad C_r = \bigcap_{n=1}^{\infty}B_{n,r}, \qquad B_{n,r} = \bigcup_{k=n}^{\infty}\Big\{\Big|\frac{S_k}{k}\Big|>\frac1r\Big\}.
\]
Достатъчно е да покажем $\mathbb P(C_r)=0$ за всяко $r\ge1$. Понеже
$(B_{n,r})_{n\ge1}$ намалява по $n$, $\mathbb P(C_r)=\lim_{n\to
\infty}\mathbb P(B_{n,r})$. По субадитивност (Тема~2):
\[
  \mathbb P(B_{n,r}) \le \sum_{k=n}^{\infty} \mathbb P\Big(\Big|\frac{S_k}{k}\Big|>\frac1r\Big).
\]
Прилагаме обобщеното неравенство на Чебишов с~$n=4$-ти момент:
\[
  \mathbb P\Big(\Big|\frac{S_k}{k}\Big|>\frac1r\Big) \le \frac{r^4}{k^4}\,\mathbb E S_k^4 = \frac{r^4}{k^4}\,\mathbb E\Big(\sum_{i=1}^k X_i\Big)^4.
\]
Разгръщаме четвъртата степен:
\[
  \Big(\sum_{i=1}^k X_i\Big)^4 = \sum_{i=1}^k X_i^4 + 6\sum_{i<j} X_i^2X_j^2 + (\text{членове с~нечетни степени}).
\]
Понеже $(X_i)$ са независими и~$\mathbb EX_i=0$, всеки член с~нечетна степен на кой да е от множителите (напр.\ $X_i^3X_j$) се
анулира при вземане на очакване: $\mathbb E[X_i^3X_j]=\mathbb
EX_i^3\cdot\mathbb EX_j = \mathbb EX_i^3\cdot0=0$. Остават само
членовете от четни степени:
\[
  \mathbb E\Big(\sum_{i=1}^k X_i\Big)^4 = k\,\mathbb EX_1^4 + 6\binom{k}{2}\big(\mathbb EX_1^2\big)^2 \le C\cdot k^2
\]
за подходяща константа $C>0$ (защото първият член расте линейно по
$k$, а~вторият - квадратично, откъдето изразът е от порядъка на
$k^2$ за големи $k$). Заместваме:
\[
  \mathbb P\Big(\Big|\frac{S_k}{k}\Big|>\frac1r\Big) \le \frac{r^4}{k^4}\cdot C k^2 = \frac{Cr^4}{k^2}.
\]
Следователно
\[
  \mathbb P(B_{n,r}) \le \sum_{k=n}^{\infty}\frac{Cr^4}{k^2} = Cr^4\sum_{k=n}^{\infty}\frac1{k^2} \xrightarrow[n\to\infty]{} 0
\]
(опашка на сходящ ред $\sum 1/k^2$ клони към нула). Следователно
$\mathbb P(C_r) = \lim_n\mathbb P(B_{n,r}) = 0$ за всяко $r\ge1$, и~по
субадитивност $\mathbb P(L^c) = \mathbb P\big(\bigcup_r C_r\big) \le
\sum_r \mathbb P(C_r) = 0$, т.е.\ $\mathbb P(L)=1$. \qedbg
\end{proofbox}

\begin{rembox}[Обобщение при \texorpdfstring{$\mathbb EX_1=\mathbb EX_i$}{EX1=EXi}]
Независимост и~общо (крайно) очакване \textbf{сами по себе си не са
достатъчни} за СЗГЧ/УЗГЧ при неeднакво разпределени величини -
нужно е допълнително условие върху дисперсиите. Едно просто
достатъчно условие е \emph{равномерна ограниченост на дисперсиите}:
съществува $M>0$, такова че $DX_i\le M$ за всяко $i$. Тогава
$\sum_{i=1}^n(X_i-\mathbb EX_i)/n = \sum_{i=1}^n X_i/n-\mathbb EX_1$,
и~заключението на СЗГЧ остава в~сила:
\[
  \frac{\sum_{i=1}^n X_i}{n}\xrightarrow{\mathbb P}\mathbb EX_1,
\]
защото по независимост и~равномерна ограниченост,
$D\big(\sum_{i=1}^n(X_i-\mathbb EX_i)\big)=\sum_{i=1}^nDX_i\le nM$, и~неравенството на Чебишов дава (аналогично на доказателството по-горе)
$\mathbb P\big(|\sum_{i=1}^n(X_i-\mathbb EX_i)/n|>\varepsilon\big)\le
M/(n\varepsilon^2)\to0$. За УЗГЧ равномерната ограниченост на
дисперсиите също е достатъчна (чрез по-финото условие на Колмогоров
$\sum_i DX_i/i^2<\infty$), но доказателството излиза извън рамките на
този курс.
\end{rembox}
